\documentclass{article}
\title{TypeScript 入門 — 型から学ぶモダン開発}
\author{TeX 図書館}

\begin{document}

\section{TypeScript とは — なぜ型が必要か}

TypeScript は Microsoft が開発した JavaScript の\textbf{スーパーセット}です。すべての JavaScript コードはそのまま TypeScript として動き、そこに\textbf{静的な型システム}を追加します。

\subsection{JavaScript の問題点}

次の JavaScript コードを見てください。

\begin{lstlisting}[language=JavaScript]
function formatPrice(price) {
  return "¥" + price.toLocaleString();
}

formatPrice(1200);        // "¥1,200" — うまくいく
formatPrice("1200");      // "¥1200"   — 意図と違うが黙って動く
formatPrice(undefined);   // 実行時エラー
\end{lstlisting}

文字列を渡してもエラーにならず、\texttt{undefined} は実行時まで気づけません。プログラムが大きくなると、この「黙って変な動きをする」性質が最大の敵になります。TypeScript では型を書くことで、こうしたミスを\textbf{コードを実行する前に}発見できます。

\begin{lstlisting}[language=TypeScript]
function formatPrice(price: number): string {
  return "¥" + price.toLocaleString();
}

formatPrice(1200);      // OK
formatPrice("1200");    // コンパイルエラー: string は number に代入できない
\end{lstlisting}

\subsection{TypeScript の動き方}

TypeScript はブラウザや Node.js がそのまま実行できる形式ではありません。\textbf{コンパイル}(トランスパイル)という変換作業で JavaScript に落とします。

\begin{itemize}
\item \texttt{tsc}(TypeScript Compiler)が型チェックと変換を行う
\item 型は実行時には\textbf{完全に消える}(型消去、type erasure)。実行速度に影響しない
\item 型チェックとコンパイル出力は独立しており、エラーがあっても JS 出力は可能(\texttt{--noEmitOnError} で制御)
\end{itemize}

つまり TypeScript の型は\textbf{開発時の安全装置であり設計の道具}です。実行環境に追加物を持ち込まない設計は、JavaScript 資産との互換性を保つ賢い選択でした。

\begin{quote}
\textbf{【演習】}\ 「TypeScript は実行を遅くする」という説の正しさを、型消去の観点から説明せよ。
\end{quote}

\textbf{解答の指針:} 誤り。型情報はコンパイル時に消えるため、実行されるのは普通の JavaScript であり、実行時の性能への影響はない(逆に、型のおかげで初期段階で最適化しやすいコード構造になる恩恵は開発者にある)。

\section{基本の型 — プリミティブとリテラル型}

\subsection{7つの基本型}

TypeScript が扱う基本の型を押さえます。

\begin{lstlisting}[language=TypeScript]
let count: number = 42;        // 数値(整数も小数も number)
let name: string = "Alice";    // 文字列
let ok: boolean = true;        // 真偽値
let ids: number[] = [1, 2, 3]; // 配列
let pair: [string, number] = ["age", 20]; // タプル(長さ固定)
let maybe: string | null = null;          // ユニオン(後述)
let cb: (n: number) => boolean = (n) => n > 0; // 関数型
\end{lstlisting}

JavaScript に馴染みがない人向けの注意点: 数値は \texttt{int} と \texttt{float} に分かれずすべて \texttt{number} です(\texttt{1} も \texttt{3.14} も同じ型)。BigInt が必要な特別なケース以外は \texttt{number} で十分です。

\subsection{リテラル型と const}

文字列や数値の\textbf{特定の値だけ}を許す型を\textbf{リテラル型}といいます。

\begin{lstlisting}[language=TypeScript]
let direction: "up" | "down" | "left" | "right";
direction = "up";    // OK
direction = "north"; // エラー: 存在しない向き

const size = 3;      // const は自動的にリテラル型 3 になる
let size2 = 3;       // let は広めの number に推論される
\end{lstlisting}

\texttt{const} で宣言した変数は値が変わらないので、TypeScript は可能な限り狭い型(リテラル型)を推論します。この「推論の狭さ」の感覚は TypeScript 力の基本です。

\subsection{型推論}

型注釈を省略しても、TypeScript は初期値から型を\textbf{推論}します。

\begin{lstlisting}[language=TypeScript]
let msg = "hello";   // string と推論される
msg = 42;            // エラー

// 引数は推論できないので注釈が必須
function double(n: number) {
  return n * 2;      // 戻り値は number と推論される
}
\end{lstlisting}

原則として「変数・戻り値は推論に任せ、\textbf{引数と公開 API には明示的に型を書く}」のが良いバランスです。

\begin{quote}
\textbf{【演習】}\ 次の変数の推論結果を予測せよ。
\begin{verbatim}
const a = [1, 2, 3];
let b = a.length > 2 ? "long" : "short";
\end{verbatim}
\end{quote}

\textbf{解答の指針:} \texttt{a} は \texttt{number[]}、\texttt{b} は \texttt{"long" | "short"}(三項演算子の両辺のリテラル型のユニオン)。ただし \texttt{b} は \texttt{let} なので、そのユニオン型に推論される。

\section{オブジェクトの型 — interface と type}

\subsection{オブジェクト型の宣言}

現実のデータはオブジェクトです。構造を型として書きます。

\begin{lstlisting}[language=TypeScript]
interface User {
  id: number;
  name: string;
  email?: string;        // ? は省略可能プロパティ
  readonly createdAt: Date; // readonly は変更不可
}

const user: User = {
  id: 1,
  name: "Alice",
  createdAt: new Date(),
};

user.name = "Bob";      // OK
user.createdAt = new Date(); // エラー: readonly
\end{lstlisting}

\texttt{email} のような省略可能プロパティは \texttt{string | undefined} を意味します。「鍵がない」ことと「値が \texttt{undefined}」の区別は TypeScript では重要です。

\subsection{interface と type の違い}

同じオブジェクト型を \texttt{type} でも書けます。

\begin{lstlisting}[language=TypeScript]
type UserId = number;
type UserAlias = {
  id: UserId;
  name: string;
};
\end{lstlisting}

使い分けの実務慣行:
\begin{itemize}
\item \textbf{オブジェクトの形状には \texttt{interface}}: 宣言マージ(同名 interface の自動統合)や extends に対応し、エラーメッセージが読みやすい
\item \textbf{ユニオン・マップ型・ユーティリティには \texttt{type}}: interface では表現できない型が書ける
\end{itemize}

\subsection{構造的部分型}

TypeScript の型は\textbf{名前}ではなく\textbf{構造}で判定されます( structural typing 、ダックタイピングの静的版)。

\begin{lstlisting}[language=TypeScript]
interface Point { x: number; y: number; }
interface Pos3D { x: number; y: number; z: number; }

const p: Point = { x: 1, y: 2, z: 3 };
// OK! Pos3D は Point が要求する性質をすべてもつ
\end{lstlisting}

\texttt{Pos3D} が \texttt{Point} を継承していなくても、必要なプロパティをもつので代入できます。「必要なものが揃っているか」だけを見るのが構造的部分型です。この規則のおかげで、異なるライブラリ間の型でも自然に互換します。

\begin{quote}
\textbf{【演習】}\ 次のコードでエラーになる行をすべて見つけよ。
\begin{verbatim}
interface Animal { name: string; }
interface Dog { name: string; bark(): void; }

const d: Dog = { name: "Pochi", bark() {} };
const a: Animal = d;          // (1)
const d2: Dog = { name: "X" }; // (2)
\end{verbatim}
\end{quote}

\textbf{解答の指針:} (1) は OK — Dog は Animal の要求(name)を満たす。(2) はエラー — bark が足りない。

\section{関数の型 — 引数、戻り値、オーバーロード}

\subsection{関数型の基本}

\begin{lstlisting}[language=TypeScript]
// 引数に型、戻り値の型は推論(明示も可)
function sum(...nums: number[]): number {
  return nums.reduce((a, b) => a + b, 0);
}

// アロー関数でも同じ
const multiply = (a: number, b: number): number => a * b;

// デフォルト引数・省略可能引数
function greet(name: string, greeting = "Hello"): string {
  return `${greeting}, ${name}`;
}
\end{lstlisting}

\subsection{void と never}

\texttt{void} は「戻り値がない」ことを表し、\texttt{never} は「戻ってこない」ことを表します。

\begin{lstlisting}[language=TypeScript]
function log(msg: string): void {
  console.log(msg);   // 値を返さない
}

function fail(message: string): never {
  throw new Error(message);  // 必ず例外で止まる
}

function infiniteLoop(): never {
  while (true) {}
}
\end{lstlisting}

\texttt{never} は到達不可能な値の型で、後で学ぶ「網羅性チェック」で重要な役割を果たします。

\subsection{コールバックと関数型の再利用}

\begin{lstlisting}[language=TypeScript]
type Comparator<T> = (a: T, b: T) => number;

function sortStringsBy(cmp: Comparator<string>) {
  // ...
}

const byLength: Comparator<string> = (a, b) => a.length - b.length;
\end{lstlisting}

コールバックを受け取る高階関数は JavaScript/TypeScript の日常です。\texttt{Array.map} や \texttt{filter} の型をエディタで確認すると、ジェネリクス(第8節)で書かれていることが分かります。

\begin{quote}
\textbf{【演習】}\ 次の関数に型を付けよ。「商品名と単価の配列を受け取り、消費税込み価格の一覧を返す」。
\begin{verbatim}
function withTax(items, rate = 0.1) {
  return items.map(([name, price]) => [name, price * (1 + rate)]);
}
\end{verbatim}
\end{quote}

\textbf{解答の指針:} \texttt{function withTax(items: [string, number][], rate = 0.1): [string, number][]}。

\section{ユニオン型と型の絞り込み}

\subsection{ユニオンで「いずれか」を表す}

状態が複数ありえる値は、\textbf{ユニオン型}で表現します。

\begin{lstlisting}[language=TypeScript]
type Result =
  | { status: "success"; data: string[] }
  | { status: "error"; message: string };

function show(r: Result): string {
  if (r.status === "success") {
    return r.data.join(", ");  // ここでは r が成功側だと分かる
  }
  return `Error: ${r.message}`; // ここではエラー側だと分かる
}
\end{lstlisting}

\texttt{if} 文で \texttt{r.status} を調べると、各分岐の中で TypeScript は \texttt{r} の型を\textbf{絞り込んで}くれます(型の絞り込み、narrowing)。判別子(\texttt{status})をもつユニオンは\textbf{判別可能ユニオン}と呼ばれ、現実的な TypeScript 設計の最重要パターンです。

\subsection{typeof と in による絞り込み}

\begin{lstlisting}[language=TypeScript]
function process(value: string | number): number {
  if (typeof value === "string") {
    return value.length;  // string として扱える
  }
  return value * 2;       // number として扱える
}

function area(shape: { kind: "circle"; r: number } | { kind: "rect"; w: number; h: number }) {
  if ("r" in shape) return Math.PI * shape.r ** 2;
  return shape.w * shape.h;
}
\end{lstlisting}

\subsection{never による網羅性チェック}

判別可能ユニオンに新しいケースを追加したとき、処理側の書き漏れをコンパイルエラーで検出できます。

\begin{lstlisting}[language=TypeScript]
type Shape =
  | { kind: "circle"; r: number }
  | { kind: "rect"; w: number; h: number }
  | { kind: "triangle"; base: number; h: number };

function area(s: Shape): number {
  switch (s.kind) {
    case "circle": return Math.PI * s.r ** 2;
    case "rect": return s.w * s.h;
    case "triangle": return (s.base * s.h) / 2;
    default: {
      const _exhaustive: never = s; // 全ケース処理済みなら never
      return _exhaustive;
    }
  }
}
\end{lstlisting}

\texttt{Shape} に \texttt{"ellipse"} を追加すると、\texttt{default} 節の \texttt{never} 代入がエラーになり「処理漏れ」を教えてくれます。これが \texttt{never} 型の実務的価値です。

\begin{quote}
\textbf{【演習】}\ API の応答を表す型 \texttt{Response = \{ ok: true; body: string \} | \{ ok: false; code: number \}} を判別可能ユニオンとして書き、両方のケースを処理する関数を \texttt{never} による網羅チェックつきで実装せよ。
\end{quote}

\textbf{解答の指針:} 判別子は \texttt{ok}。\texttt{if (r.ok)} で成功側、\texttt{else} でエラー側を絞り込む。switch 文 + \texttt{default: const \_x: never = r} でもよい。

\section{ジェネリクス — 型を引数にする}

\subsection{型パラメータ}

「型を引数として受け取る関数や型」を\textbf{ジェネリクス}といいます。重複コードを型安全に再利用するための道具です。

\begin{lstlisting}[language=TypeScript]
function first<T>(items: T[]): T | undefined {
  return items[0];
}

first([1, 2, 3]);        // T = number、戻り値 number | undefined
first(["a", "b"]);       // T = string
first<number>([1, 2]);   // 明示指定もできる
\end{lstlisting}

\texttt{T} は「呼び出し時に決まる型」の名札です。配列要素の型が何であっても、\textbf{入力と出力の型の対応}を保てるのがジェネリクスの価値です。

\subsection{制約付きジェネリクス}

\texttt{extends} で「T は少なくともこれを満たす」という制約を付けられます。

\begin{lstlisting}[language=TypeScript]
// length をもつものだけを受け取る
function totalLength<T extends { length: number }>(a: T, b: T): number {
  return a.length + b.length;
}

totalLength("ab", "cd");   // OK: string は length をもつ
totalLength([1], [2, 3]);  // OK: 配列も length をもつ
totalLength(3, 4);         // エラー: number に length はない
\end{lstlisting}

\subsection{標準ライブラリのジェネリクス}

\begin{lstlisting}[language=TypeScript]
const map = new Map<string, number>();
map.set("alice", 90);

const nums: Array<number> = [1, 2, 3];
const found: number | undefined = nums.find((n) => n > 1);

// Promise もジェネリクス: Promise<T> は「T に解決する非同期計算」
async function fetchUser(id: number): Promise<User> {
  const res = await fetch(`/api/users/${id}`);
  return res.json() as Promise<User>;
}
\end{lstlisting}

\begin{quote}
\textbf{【演習】}\ 次の関数をジェネリクスで型付けせよ。「キーでオブジェクトの配列をグループ化する」。
\begin{verbatim}
function groupBy(items, key) {
  const result = {};
  for (const item of items) {
    const k = item[key];
    (result[k] ??= []).push(item);
  }
  return result;
}
\end{verbatim}
\end{quote}

\textbf{解答の指針:} \texttt{function groupBy<T, K extends keyof T>(items: T[], key: K): Record<string, T[]>} の形が理想(実装詳細はキーの型に依存するので、簡易に \texttt{Record<string, T[]>} でよい)。

\section{クラスとオブジェクト指向}

\subsection{クラスの基本}

TypeScript のクラスは、フィールド宣言・アクセス修飾子・実装隠蔽を型レベルで支援します。

\begin{lstlisting}[language=TypeScript]
class BankAccount {
  private balance: number = 0;   // 外部から触れない

  constructor(public readonly owner: string) {}

  deposit(amount: number): void {
    if (amount <= 0) throw new Error("金額は正の数");
    this.balance += amount;
  }

  getBalance(): number {
    return this.balance;
  }
}

const acc = new BankAccount("Alice");
acc.deposit(1000);
acc.balance = -9999; // エラー: private
\end{lstlisting}

\subsection{継承とインターフェース実装}

\begin{lstlisting}[language=TypeScript]
interface Shape {
  area(): number;
}

class Circle implements Shape {
  constructor(private r: number) {}
  area(): number { return Math.PI * this.r ** 2; }
}

class Square implements Shape {
  constructor(private side: number) {}
  area(): number { return this.side ** 2; }
}

const shapes: Shape[] = [new Circle(1), new Square(2)];
const total = shapes.reduce((sum, s) => sum + s.area(), 0);
\end{lstlisting}

\texttt{implements} は「この interface を満たすことを約束する」宣言です。配列 \texttt{shapes} の要素型が \texttt{Shape} なので、\texttt{area()} をもつことだけを仮定して処理でき、新しい図形クラスを追加してもこのコードは無変更で動きます。\textbf{多態}(ポリモーフィズム)の実務的価値です。

\subsection{public / private / protected}

\begin{itemize}
\item \textbf{public}(既定): どこからでもアクセス可
\item \textbf{private}: そのクラスの内部からのみ
\item \textbf{protected}: そのクラスとサブクラスから
\end{itemize}

これらの修飾子はコンパイル時のみ有効で、実行される JS には残りません。実行時に強固に隠したい場合は ES の \texttt{\#private} フィールドを使います。

\begin{quote}
\textbf{【演習】}\ \texttt{class Temperature} を作れ。摂氏を内部に保持し、\texttt{celsius} の getter/setter(負273.15度未満はエラー)、\texttt{toFahrenheit()} メソッドをもたせよ。
\end{quote}

\textbf{解答の指針:} \texttt{private c: number} を保持、setter で \texttt{if (v < -273.15) throw new Error(...)}. \texttt{toFahrenheit() \{ return this.c * 9 / 5 + 32; \}}。

\section{非同期処理とエラーハンドリング}

\subsection{Promise と async/await}

非同期処理は \texttt{Promise<T>}、「将来 \texttt{T} を受け渡す約束」として型付けされます。

\begin{lstlisting}[language=TypeScript]
function delay(ms: number): Promise<void> {
  return new Promise((resolve) => setTimeout(resolve, ms));
}

async function loadAll(ids: number[]): Promise<User[]> {
  const users: User[] = [];
  for (const id of ids) {
    await delay(100);
    users.push({ id, name: `user${id}`, createdAt: new Date() });
  }
  return users;
}

// 並列実行: Promise.all は「成功なら T[]、失敗なら即エラー」
const all = await Promise.all(ids.map((id) => fetchUser(id)));
\end{lstlisting}

\texttt{await} は Promise の中身を取り出す操作で、\texttt{async} 関数の中でのみ使えます。\texttt{async} 関数の戻り値は必ず Promise になる点に注意しましょう。

\subsection{エラーハンドリング}

\begin{lstlisting}[language=TypeScript]
// Result 風の判別可能ユニオンで「失敗しうる関数」を型で表現
type Result<T> = { ok: true; value: T } | { ok: false; error: string };

function parseAge(input: string): Result<number> {
  const n = Number(input);
  if (Number.isNaN(n) || n < 0) {
    return { ok: false, error: "0以上の数を入力してください" };
  }
  return { ok: true, value: n };
}

const r = parseAge("25");
if (r.ok) {
  console.log(r.value);  // ここでは number
} else {
  console.error(r.error); // ここでは string
}
\end{lstlisting}

例外(\texttt{throw})は TypeScript でも使えますが、\texttt{catch} された値は \texttt{unknown} 型になります。そのため現代の TypeScript では「失敗の可能性を戻り値の型に表す」判別可能ユニオン方式が好まれる場面が増えています。

\begin{quote}
\textbf{【演習】}\ \texttt{async function fetchJson<T>(url: string): Promise<T>} を、\texttt{fetch} が失敗した場合と \texttt{res.ok} が false の場合を区別するエラーメッセージつきで実装せよ。
\end{quote}

\textbf{解答の指針:} \texttt{try} ブロックで \texttt{const res = await fetch(url)}、\texttt{if (!res.ok) throw new Error(\`HTTP \${res.status}\`)}、\texttt{return await res.json() as T}。

\section{実践演習 — 型でリファクタリングする}

この節では、「動くけれど型が弱い」コードを 3 題リファクタリングします。各題とも「まず問題点を自分の言葉で説明する」のが正しい手順です。型は最後の道具です。

\subsection{第 1 題 — 眠っているユニオン}

\begin{lstlisting}[language=TypeScript]
// 変更前: kind はただの文字列で、中身の整合性は実行時次第
function render(node: { kind: string; text?: string; src?: string }) {
  if (node.kind === "text") return node.text ?? "";
  if (node.kind === "image") return `<img src="${node.src}">`;
  return "";
}
\end{lstlisting}

\textbf{解答の指針:}\ 問題点は (1) \texttt{kind} にどんな文字列でも入る、(2) \texttt{text} と \texttt{src} が常に省略可能で「text ノードなのに src がある」不正状態が作れる、(3) 新しい kind を追加しても処理漏れが検出されない。判別可能ユニオン + \texttt{never} による網羅チェックで解決する:
\begin{lstlisting}[language=TypeScript]
type Node =
  | { kind: "text"; text: string }
  | { kind: "image"; src: string };

function render(node: Node): string {
  switch (node.kind) {
    case "text": return node.text;
    case "image": return `<img src="${node.src}">`;
    default: { const _x: never = node; return _x; }
  }
}
\end{lstlisting}

\subsection{第 2 題 — any が広げる穴}

\begin{lstlisting}[language=TypeScript]
// 変更前: 入力をそのまま any で受けている
function pick(objs: any[], key: any): any[] {
  return objs.map((o) => o[key]);
}
\end{lstlisting}

\textbf{解答の指針:}\ \texttt{any} はコンパイラの目を止めるので、\texttt{pick(users, "nmae")}(typo)も通ってしまう。ジェネリクスと \texttt{keyof} で「存在するキーだけ」に限定する:
\begin{lstlisting}[language=TypeScript]
function pick<T, K extends keyof T>(objs: T[], key: K): T[K][] {
  return objs.map((o) => o[key]);
}
// pick(users, "nmae") はコンパイルエラーになる
\end{lstlisting}

\subsection{第 3 題 — 状態の組み合わせ}

\begin{lstlisting}[language=TypeScript]
// 変更前: boolean と optional の組み合わせで不正状態を許す
interface Panel {
  isLoading: boolean;
  error?: string;
  data?: User[];
}
\end{lstlisting}

\textbf{解答の指針:}\ 「ローディング中なのにエラーとデータがある」など、実行時の不変条件が型に書かれていない。第 7 節の判別可能ユニオンで状態機械にする:
\begin{lstlisting}[language=TypeScript]
type Panel =
  | { state: "loading" }
  | { state: "error"; message: string }
  | { state: "ready"; data: User[] };
\end{lstlisting}
描画側の \texttt{switch} は網羅チェックつきで書くため、状態の追加(\texttt{"empty"} など)に対してコンパイラが処理漏れを指摘してくれます。\textbf{「ありえない状態を型から消す」ことがリファクタリングの核心}です。

\begin{quote}
\textbf{【演習】}\ 第 1 題のコードに \texttt{\{ kind: "button"; label: string \}} を追加したとき、変更前と変更後のそれぞれで起きることを比較して説明せよ。
\end{quote}

\textbf{解答の指針:}\ 変更前は \texttt{kind} が \texttt{string} 型なので、render の if 文に case を足し忘れても黙って空文字を返す(バグは画面上で発覚)。変更後は \texttt{never} 代入がエラーになり、コンパイルの時点で追加箇所が分かる。

\section{ユーティリティ型と実践設計}

\subsection{組み込みユーティリティ型}

既存の型から別の型を\textbf{変換する}組み込み型たちを紹介します。

\begin{lstlisting}[language=TypeScript]
interface User {
  id: number;
  name: string;
  email: string;
}

type PartialUser = Partial<User>;      // 全プロパティ省略可能
type RequiredUser = Required<User>;    // 全プロパティ必須
type UserPreview = Pick<User, "id" | "name">; // 部分取り出し
type UserUpdate = Omit<User, "id">;    // id を除く
type NameOrEmail = "name" | "email";   // keyof User と同じ
type UserMap = Record<string, User>;   // キー -> User の辞書
\end{lstlisting}

特に \texttt{Partial}(更新用 DTO)、\texttt{Pick}/\texttt{Omit}(画面ごとに必要な形だけを切り出す)は実務で毎日のように登場します。\textbf{元の interface を一箇所で定義し、派生型は変換で作る}のが保守しやすい設計です。

\subsection{keyof と typeof}

型と値の世界をつなぐ演算子です。

\begin{lstlisting}[language=TypeScript]
const config = {
  host: "localhost",
  port: 8080,
  debug: true,
};

type Config = typeof config;             // 値から型を作る
type ConfigKey = keyof Config;           // "host" | "port" | "debug"

function get<K extends keyof Config>(key: K): Config[K] {
  return config[key];
}

const port = get("port");  // number 型(正確に推論される)
\end{lstlisting}

\texttt{get("port")} の戻り値が単なる \texttt{string | number | boolean} でなく \texttt{number} になるのがポイントです。\texttt{K extends keyof Config} と \texttt{Config[K]}(ルックアップ型)の組み合わせで、キーと値の型の対応を保てます。

\subsection{実践:型安全な設計の原則}

最後に、この教科書全体を貫く設計原則をまとめます。

\begin{enumerate}
\item \textbf{不正な状態を表現できないようにする}: \texttt{string} 1つで状態を持たず、判別可能ユニオンやリテラル型の組合せで「ありえない値」を型レベルで排除する
\item \textbf{境界で検証し、内部では信頼する}: API 応答や入力は \texttt{unknown} として受け取り、検証を通してから具体型にする
\item \textbf{推論に任せられる場所は任せる}: 型を書きすぎると可読性が落ちる。引数・公開境界に注釈を集中させる
\item \textbf{派生型は変換で作る}: \texttt{Pick}・\texttt{Omit}・\texttt{Record} で元の型から作り、手書きの重複をなくす
\end{enumerate}

\begin{quote}
\textbf{【総合演習】}
\begin{enumerate}
\item Todo 管理アプリの状態を、\texttt{Todo \{ id: number; title: string; done: boolean \}} から出発して「未保存 / 保存済み / 削除済み」を区別する判別可能ユニオンで設計せよ。
\item \texttt{Todo[]} から未完了の \texttt{Todo} だけを抽出する関数を、\texttt{Omit} を使って「title だけのリスト」も作れるよう汎用化せよ。
\item 「不正な状態を表現できないようにする」原則に従い、\texttt{isLoading: boolean} と \texttt{data?: User} を別々にもつ設計の問題点を指摘せよ。
\end{enumerate}
\end{quote}

\textbf{解答の指針:} 1. \texttt{| \{ state: "draft"; todo: Todo \} | \{ state: "saved"; todo: Todo \} | \{ state: "deleted"; id: number \}} の形。削除済みに todo 本文は不要なので \texttt{id} だけをもたせると表現が正確。2. \texttt{function pluck<T, K extends keyof T>(items: T[], key: K): T[K][] \{ return items.map((it) => it[key]); \}}。3. \texttt{isLoading: true} なのに \texttt{data} も存在する、両方 \texttt{false/undefined} の状態が起きうる——「ありえない組合せ」を型が許してしまう。\texttt{| \{ kind: "loading" \} | \{ kind: "ready"; data: User \} | \{ kind: "error"; message: string \}} とすべき。

\end{document}
